\documentclass[11pt]{article}
\usepackage[utf8]{inputenc}
\usepackage{amsmath,amssymb}
\usepackage{hyperref}
\title{Robust Genome Wide Association Study under Distribution Shift}
\author{Anonymous}
\date{}
\begin{document}
\maketitle

\begin{abstract}
This paper studies Genome Wide Association Study. Grounded in a real, citable literature \cite{ref1,ref2,ref3,ref4,ref5,ref6,ref7,ref8},
we propose a focused method and report \textbf{illustrative} experiments.
All references below are real and verifiable (OpenAlex/arXiv); the reported
numbers are simulated to illustrate intended behaviour.
\end{abstract}

\section{Introduction}
Genome Wide Association Study is an active research area. This work targets a concrete gap identified from
the recent literature and contributes a method together with an evaluation protocol.

\section{Related Work (real literature)}
The following works, retrieved from public bibliographic APIs, frame our study:
\begin{itemize}
\item Dunham et al. (2012) — "An integrated encyclopedia of DNA elements in the human genome", Nature (cited 19581).
\item Chang et al. (2015) — "Second-generation PLINK: rising to the challenge of larger and richer datasets", GigaScience (cited 14199).
\item Price et al. (2006) — "Principal components analysis corrects for stratification in genome-wide association studies", Nature Genetics (cited 10822).
\item Bowden et al. (2015) — "Mendelian randomization with invalid instruments: effect estimation and bias detection through Egger regression", International Journal of Epidemiology (cited 10750).
\item Prior work (2013) — "The Genotype-Tissue Expression (GTEx) project.", PubMed (cited 10115).
\item Bycroft et al. (2018) — "The UK Biobank resource with deep phenotyping and genomic data", Nature (cited 9924).
\end{itemize}

\section{Method}
We model distribution shift explicitly and optimise a worst-case objective over an uncertainty set of plausible test distributions. We instantiate this idea for Genome Wide Association Study and analyse its cost/quality trade-off.

\section{Experiments (Illustrative)}
\begin{center}
\begin{tabular}{lcc}
System & Primary Metric & Efficiency \\ \hline
Strong baseline & 0.812 & 1.00$\times$ \\
Ablation & 0.828 & 0.92$\times$ \\
\textbf{Ours} & \textbf{0.851} & \textbf{0.71$\times$}
\end{tabular}
\end{center}
\emph{Metrics simulated to illustrate the intended effect; not measured.}

\section{Conclusion}
We connected a real literature to a concrete method for Genome Wide Association Study. Future work will
validate the illustrative results with real experiments.

\begin{thebibliography}{9}
\bibitem{ref1} Ian Dunham, Ewan Birney, Javier Herrero et al.. An integrated encyclopedia of DNA elements in the human genome. \emph{Nature}, 2012. DOI:10.1038/nature11247
\bibitem{ref2} Christopher Chang, Carson C. Chow, Laurent CAM Tellier et al.. Second-generation PLINK: rising to the challenge of larger and richer datasets. \emph{GigaScience}, 2015. DOI:10.1186/s13742-015-0047-8
\bibitem{ref3} Alkes L. Price, Nick J. Patterson, Robert M. Plenge et al.. Principal components analysis corrects for stratification in genome-wide association studies. \emph{Nature Genetics}, 2006. DOI:10.1038/ng1847
\bibitem{ref4} Jack Bowden, George Davey Smith, Stephen Burgess. Mendelian randomization with invalid instruments: effect estimation and bias detection through Egger regression. \emph{International Journal of Epidemiology}, 2015. DOI:10.1093/ije/dyv080
\bibitem{ref5} . The Genotype-Tissue Expression (GTEx) project.. \emph{PubMed}, 2013. DOI:10.1038/ng.2653
\bibitem{ref6} Clare Bycroft, Colin Freeman, Desislava Petkova et al.. The UK Biobank resource with deep phenotyping and genomic data. \emph{Nature}, 2018. DOI:10.1038/s41586-018-0579-z
\bibitem{ref7} Jack Bowden, George Davey Smith, Philip Haycock et al.. Consistent Estimation in Mendelian Randomization with Some Invalid Instruments Using a Weighted Median Estimator. \emph{Genetic Epidemiology}, 2016. DOI:10.1002/gepi.21965
\bibitem{ref8} Paul R. Burton, David G. Clayton, Lon R. Cardon et al.. Genome-wide association study of 14,000 cases of seven common diseases and 3,000 shared controls. \emph{Nature}, 2007. DOI:10.1038/nature05911
\end{thebibliography}
\end{document}
